% \section{提案システムの設計方針}
\section{提案システムの要件}
本システムは，帰宅時の偶発的な対面コミュニケーションの促進を目的とし，その実現にあたりユーザの自由意思の尊重，心理的負担の低減，および個人のプライバシーの保護を設計の軸としている．
具体的には，
\begin{itemize}
  \item 提示情報の確認過程や投票において，ユーザの意図や関心の有無が他者に推測されない
  \item 特定の人物に対してのみ情報が共有される状況を避ける
  \item 帰宅を選択しない場合や投票を行わない場合でも不自然にならない
  \item 各ユーザの帰宅予測時刻や滞在状況など，個人の行動を特定可能な情報を公開しない
\end{itemize}
の4つの要件を満たす設計を目指した．


\paragraph{ユーザの意図や関心が推測されない}
提示情報の確認状況や投票行動が他者に推測可能な状況では，ユーザは周囲の目を意識せざるを得ず本来の意思とは異なる選択をしてしまう可能性がある．
例えば，誰がどの時刻に投票したかという情報が可視化されると，同行したくない相手が投票していた場合に断りにくくなったり，一度投票した時刻に必ず帰宅しなければならないという義務感が生じたりするおそれがある．
また，提示情報の確認に何らかの操作や行動が伴う場合，投票数が少ない状況ではその行動から投票者であると推測されてしまう可能性がある．
このような状況は，ユーザに心理的圧力を与え，自由な選択を阻害する要因となる．
そこで本システムでは，提示情報の確認過程や投票行動まで，ユーザの意図や関心の有無が他者に推測されないよう配慮した設計とした．


\paragraph{特定の人物に対してのみ情報が共有されない}
情報が特定の人物に対してのみ提示される状況では，その人物が暗黙的に行動を期待されていると受け取りやすく，同行しないという選択を取りにくくなる可能性がある．
例えば，「AさんとBさんが同じ時刻に帰れそうである」といった形で特定の人物の組み合わせを提示した場合，その相手との同行を前提とした期待が生じやすく，同行しない場合に断る理由を考えなければならないなど，心理的負担が生じうる．
このような名指し的・指名的な情報共有は，ユーザに対して暗黙の期待や同調圧力を生みやすく，意思決定の自由度を低下させる要因となる．
そこで本システムでは，特定の人物に対してのみ情報が共有される状況を避け，全体に向けた形で情報を提示する設計とした．
これにより，特定の相手と行動を共にしたくない場合であっても，ユーザは理由の後付けが可能であり，対人関係上の心理的負担を抑えられる．


\paragraph{投票や帰宅を選択しない場合でも不自然にならない}
ユーザが投票を行わない，あるいは提示された時刻に合わせて帰宅しないといった選択をしても不自然にならないよう配慮する点も本システムにおける重要な要件である．
もし，投票や同行が事実上の前提となってしまうと，ユーザはそれらの行動を「求められている」と感じ，義務感や心理的負担を抱くおそれがある．
そこで本システムでは，ユーザの行動を直接的に誘導・強制するのではなく，あくまで選択肢としての提示を基本方針とした．
ユーザは，提示される帰宅時刻の候補に投票するかどうか，また最多投票の時刻に合わせて帰宅するかどうかを自身の判断で決定できる．
投票の有無や実際の行動は他者に公開されず，これにより不参加の選択が正当なものとして成立する設計としている．


\paragraph{個人の行動を特定可能な情報を公開しない}
個々のユーザの帰宅予測時刻や滞在状況といった情報は，個人の行動パターンの推定や追跡につながる可能性がある．
これらの情報が外部から容易に取得可能な状態にある場合，意図せず個人の行動履歴が可視化され，プライバシー侵害のリスクが生じる．
そこで本システムでは，各ユーザの帰宅予測時刻や滞在状況など，個人の行動を特定可能な情報を公開しない設計とした．
これにより，ユーザは自身の行動の追跡を過度に意識する必要がなく，安心してシステムを利用できる環境を実現している．


